% !TEX program = xelatex
% !TEX encoding = UTF-8

\documentclass[12pt, a4paper]{article}
\usepackage[UTF8]{ctex}
\usepackage{amsmath}
\usepackage{amssymb}
\usepackage{geometry}
\usepackage{enumitem}
\usepackage{fancyhdr}
\usepackage{titlesec}
\usepackage{xcolor}
\usepackage{listings}
\usepackage{framed}
\usepackage{tikz}

% 页面设置
\geometry{left=2.5cm, right=2.5cm, top=2.5cm, bottom=2.5cm}

% 页眉页脚
\pagestyle{fancy}
\fancyhf{}
\rhead{专升本数据结构讲义}
\lhead{章节精讲}
\cfoot{\thepage}

% 标题格式
\titleformat{\section}{\large\bfseries}{\thesection}{1em}{}
\titleformat{\subsection}{\normalsize\bfseries}{\thesubsection}{1em}{}

% 蓝色边框环境
\newenvironment{bluebox}{\begin{framed}\begin{minipage}{\linewidth}}{\end{minipage}\end{framed}}

% 代码样式
\lstset{
    backgroundcolor=\color{lightgray!20},
    basicstyle=\ttfamily\small,
    numbers=left,
    numberstyle=\tiny,
    frame=single,
    breaklines=true
}

\title{\textbf{第二章 线性表}}
\author{}
\date{}

\begin{document}

\maketitle

\section*{本章知识脉络}

线性表是最基本、最常用的数据结构，它是有序的元素序列。从本章开始，我们将深入学习各种具体的数据结构。线性表是后续学习栈、队列、串等数据结构的基础，也是算法面试中的重点考查内容。本章的核心在于理解顺序存储和链式存储两种实现方式的原理、优缺点，以及在不同场景下的选择依据。考生需要熟练掌握两种存储结构的操作算法及其时间空间复杂度分析。

线性表作为一种基础数据结构，广泛应用于各种实际场景。在实际软件开发中，线性表的实现随处可见：数组是最典型的顺序存储线性表，链表在各种编程语言中都有广泛使用。理解线性表的原理对于提高编程能力和解决实际问题都具有重要意义。

\section{线性表的定义与基本概念}

线性表（Linear List）是最基本的数据结构之一，它是有序的元素有限序列。线性表中元素之间存在一对一的关系，即每个元素（除第一个和最后一个外）都有且仅有一个直接前驱和一个直接后继。线性表是最简单也是最直观的数据结构，它反映了客观世界中线性关系（如排队、编号等）的抽象。

线性表的形式化定义为：线性表是n（n≥0）个同类型数据元素的有限序列。n称为线性表的长度，当n=0时，线性表为空表。线性表的第一个元素称为表头，最后一个元素称为表尾。线性表是一种线性结构，这是它与树形结构、图状结构的根本区别。

线性表有两个基本特征：一是所有元素属于同一数据对象，即线性表中的元素具有相同的数据类型；二是元素之间存在有序的先后次序关系，这种关系是线性表逻辑结构的核心。线性表的长度是动态变化的，可以在插入操作后变长，在删除操作后变短。

线性表的基本运算包括：初始化线性表、销毁线性表、判断线性表是否为空、求线性表长度、取线性表第i个元素、定位查找、插入元素、删除元素、遍历线性表等。这些基本运算是实现更复杂算法的基础，考生需要熟练掌握各种运算的算法实现和复杂度分析。

从存储方式来看，线性表可以分为顺序存储和链式存储两种基本方式。顺序存储的线性表称为顺序表，链式存储的线性表称为链表。两种存储方式各有优缺点，适用于不同的应用场景。顺序表的优点是访问效率高，缺点是插入删除需要移动元素；链表的优点是插入删除灵活，缺点是访问效率低。

\section{顺序存储结构}

顺序存储是指使用一段连续的存储单元依次存储线性表中的各个元素，使得逻辑上相邻的元素在物理位置上也相邻。顺序表是最常见的线性表存储方式，它利用数组来实现，数组的下标即对应线性表的元素位置。

顺序表的特点是：存储密度高，即一个顺序表结点占用的存储空间等于该结点数据元素本身占用的空间，没有额外的指针开销；可以实现随机访问，即可以在O(1)时间内访问顺序表中的任意元素，这是顺序表最重要的优点；插入和删除操作需要移动元素，时间复杂度为O(n)。

顺序表的结构定义通常包括一个数组和一个整型变量。数组用于存储数据元素，整型变量用于记录当前顺序表的长度。在C语言中，可以用静态数组或动态分配的一维数组来实现顺序表。动态分配的数组可以更好地利用内存空间，避免固定大小带来的限制。

顺序表的插入操作需要注意以下几点：插入位置可以在表头、表尾或表中任意位置；插入前需要判断顺序表是否已满；插入位置之后的元素需要依次向后移动一位；插入后需要将新元素放到正确位置，并更新表长。插入操作的时间复杂度为O(n)，主要时间消耗在移动元素上。

顺序表的删除操作同样需要注意：删除位置可以是表头、表尾或表中任意位置；删除前需要判断顺序表是否为空；删除位置之后的元素需要依次向前移动一位；删除后需要更新表长。删除操作的时间复杂度为O(n)，主要时间消耗在移动元素上。

顺序表的按位置查找操作可以在O(1)时间内完成，因为数组支持随机访问。按值查找操作需要遍历顺序表，平均时间复杂度为O(n)。在有序顺序表上，按值查找可以使用折半查找，时间复杂度为O(log n)，这是后续章节将要学习的内容。

顺序存储结构适用于以下场景：需要频繁访问元素而较少进行插入删除操作的情况；预先知道数据规模，可以合理分配存储空间的情况；需要实现随机访问的情况。在这些场景下，顺序表的效率优势能够得到充分发挥。

\section{链式存储结构}

链式存储是线性表的另一种重要存储方式，它使用一组任意的存储单元（不必连续）来存储线性表中的元素。链式存储通过指针显式地表示元素之间的逻辑关系，每个结点包含数据域和指针域。数据域存储数据元素本身，指针域存储指向其他结点的指针。

与顺序存储相比，链式存储的最大优点是插入和删除操作不需要移动元素，只需要修改相关结点的指针即可。链式存储不需要预先分配连续的空间，可以根据需要动态分配内存。但是，链式存储失去了随机访问的能力，访问第i个元素必须从头开始遍历，时间复杂度为O(n)。

单链表是最简单的链表结构，每个结点包含一个数据域和一个指向后继结点的指针。单链表的表尾结点的指针指向NULL，表示链表的结束。单链表可以分为带头结点的单链表和不带头结点的单链表。头结点是链表中的一个特殊结点，它不存储实际数据，仅作为链表的入口点。

头结点的作用是简化链表操作。带头结点的单链表在表头插入和删除操作时，不需要特殊处理，因为始终存在一个头结点作为链表的第一个结点。头结点使得链表的操作更加统一和简洁，减少了边界条件的处理复杂度。

单链表的基本操作包括：创建链表（头插法、尾插法）、插入结点（在指定位置插入、在表头插入、在表尾插入）、删除结点（删除指定位置结点、删除特定值结点）、查找结点（按位置查找、按值查找）、遍历链表等。每个操作都需要注意指针的修改顺序，以避免链表断裂或内存泄漏。

双链表相比单链表增加了指向前驱结点的指针，可以在O(1)时间内找到前驱结点。双链表的插入和删除操作需要同时修改前驱和后继两个方向的指针。双链表适用于需要频繁查找前驱结点的场景，但需要占用更多的存储空间。

循环链表是另一种重要的链表形式，它的最后一个结点的指针指向头结点（对于循环单链表）或指向头结点（对于循环双链表）。循环链表形成一个环，可以从任意结点出发遍历整个链表。循环单链表在需要经常在表尾插入元素的场景下很有用，配合尾指针可以实现在O(1)时间内完成表尾插入操作。

\section{顺序表与链表的比较}

顺序表和链表是线性表的两种基本存储方式，各有优缺点，适用于不同的应用场景。理解它们的区别和适用条件是本章的重点内容，也是数据结构课程的核心知识点之一。在实际应用中，需要根据具体问题的需求选择合适的存储方式。

从时间复杂度角度分析：按位置访问操作，顺序表为O(1)，链表为O(n)；按位置插入删除操作，顺序表为O(n)，链表为O(1)（前提是已经定位到位置）；按值查找操作，在有序情况下顺序表可以用折半查找为O(log n)，链表始终为O(n)。因此，访问操作频繁时选择顺序表，插入删除操作频繁时选择链表。

从空间复杂度角度分析：顺序表需要预先分配连续的空间，空间利用率高；链表每个结点需要额外的指针空间，空间利用率较低。但是，顺序表需要预先估计数据规模，可能造成空间浪费或不足；链表可以动态分配空间，更加灵活。

从存储密度角度分析：顺序表的存储密度为1，即每个存储单元都用于存储数据；链表的存储密度小于1，因为每个结点除了存储数据外还需要存储指针（单链表每个结点需要额外一个指针，双链表需要额外两个指针）。存储密度反映了存储空间的利用效率。

从实现难度角度分析：顺序表的实现简单直观，代码容易理解和调试；链表的实现需要处理指针操作，容易出错（如指针丢失、内存泄漏等）。对于初学者来说，链表是更大的挑战，需要多练习指针操作。

综合以上分析，我们可以得出以下选择原则：如果数据规模已知且稳定，主要进行查找访问操作，选择顺序表；如果数据规模动态变化，主要进行插入删除操作，选择链表；如果需要实现随机访问，选择顺序表；如果需要频繁在任意位置插入删除，选择链表。

\section{链表操作详解}

\subsection{单链表的插入操作}

单链表插入操作分为三种情况：在表头插入、在表尾插入、在指定位置插入。在表头插入是最简单的情况，只需要将新结点的指针指向原头结点的下一个结点，然后将头结点的指针指向新结点即可。在表尾插入需要遍历到链表末尾，然后修改末尾结点的指针指向新结点。在指定位置插入需要先找到插入位置的前驱结点，然后修改指针将新结点插入到链表中。

单链表插入操作的核心是正确处理指针的修改顺序。对于在p结点之后插入s结点的情况，正确的操作顺序是：先修改s结点的指针指向p结点的后继结点，再修改p结点的指针指向s结点。如果顺序颠倒，在修改p->next之后将无法找到原来的后继结点，导致链表断裂。

在表头插入结点（不带头结点的单链表）的算法步骤：首先判断链表是否为空；然后创建新结点并赋值；接着将新结点的指针指向原头结点；最后更新头指针指向新结点。带头结点的单链表在表头插入更为简单，只需要将新结点插入到头结点之后即可。

在指定位置插入结点的算法步骤：首先找到插入位置的前驱结点；然后创建新结点并赋值；接着将新结点的指针指向前驱结点的后继结点；最后修改前驱结点的指针指向新结点。需要注意处理边界情况，如插入位置为表头或表尾。

\subsection{单链表的删除操作}

单链表删除操作同样分为三种情况：删除表头结点、删除表尾结点、删除指定位置结点。删除表头结点只需要将头指针移动到下一个结点即可，但需要注意释放被删除结点的内存空间。删除表尾结点需要找到倒数第二个结点，将其指针置为NULL，然后释放原表尾结点的内存。

删除指定位置结点的算法步骤：首先找到待删除结点的前驱结点；然后修改前驱结点的指针，使其指向待删除结点的后继结点；最后释放待删除结点的内存空间。需要注意处理待删除结点为第一个结点（表头结点）的情况。

有一种特殊的删除技巧，称为“交换删除法”。对于已知待删除结点p的位置，但不知道其前驱结点的情况，可以将p的后继结点的数据复制到p结点，然后删除p的后继结点。这种方法可以避免遍历链表寻找前驱结点，但会改变结点的位置。

循环单链表的删除操作需要注意判断链表是否为空。删除循环单链表的尾结点时，除了修改尾结点的指针外，还需要将新的尾结点的指针指向头结点。其他操作与非循环单链表类似。

\subsection{双链表的操作}

双链表相比单链表增加了前驱指针，这使得在O(1)时间内访问前驱结点成为可能。双链表的插入和删除操作需要同时修改前驱和后继两个方向的指针。

双链表在p结点之后插入s结点的操作步骤：首先设置s->next = p->next；然后设置s->prior = p；接着判断p是否为尾结点，如果不是则修改p->next->prior = s；最后修改p->next = s。注意操作顺序，避免指针断裂。

双链表删除p结点的操作步骤：首先判断p是否为表头结点，如果不是则修改p->prior->next = p->next；然后判断p是否为表尾结点，如果不是则修改p->next->prior = p->prior；最后释放p结点的内存空间。双链表删除操作需要修改两个方向的指针，共需要4次指针修改。

循环双链表的操作需要特别注意边界条件。判断循环双链表是否为空的条件是L->next == L（只有一个头结点时）。在循环双链表尾部插入和删除结点时，需要将新的尾结点或新的尾结点的next指针指向头结点。

\section{链表与顺序表的复杂度分析}

线性表的时间复杂度和空间复杂度分析是本章的重点内容。不同的存储结构在不同的操作上表现出不同的性能特点，理解这些特点是正确选择存储结构的依据。

顺序表按位置访问的时间复杂度为O(1)，这是因为数组支持随机访问，可以直接通过下标计算地址访问任意元素。顺序表按位置插入和删除的平均时间复杂度为O(n)，因为需要移动平均约n/2个元素。最坏情况下（插入到表头或删除表头元素）需要移动n个元素。

链表按位置访问的时间复杂度为O(n)，因为需要从头结点开始顺序遍历，直到找到第i个结点。链表按位置插入和删除的时间复杂度为O(1)（前提是已经定位到位置），因为只需要修改相关结点的指针，不需要移动元素。链表的按值查找时间复杂度为O(n)，平均需要比较n/2个元素。

从空间复杂度来看，顺序表的空间复杂度为O(n)，需要预先分配n个存储单元。链表的空间复杂度也为O(n)，但每个结点除了存储数据外还需要存储指针。单链表每个结点需要额外一个指针的空间，双链表需要额外两个指针的空间。因此，链表的实际占用空间通常大于顺序表。

顺序存储结构的存储密度为1，链式存储结构的存储密度小于1。存储密度是指数据元素本身占用的存储量与整个结构占用的存储量之比。顺序表的存储密度高，空间利用率高；链表的存储密度低，需要额外的指针空间。

在选择存储结构时，需要综合考虑各种因素。如果数据规模已知且稳定，主要进行查找访问操作，应选择顺序表。如果数据规模动态变化，主要进行插入删除操作，应选择链表。如果需要实现随机访问，应选择顺序表。如果需要频繁在任意位置插入删除，应选择链表。

\begin{bluebox}{本章核心考点总结}

	\textbf{【核心考点一：顺序表与链表的对比】}

	这是线性表章节最重要的知识点。顺序表的优点是访问效率高(O(1))、存储密度高，缺点是插入删除需要移动元素(O(n))；链表的优点是插入删除灵活(O(1))、无需预先分配空间，缺点是访问效率低(O(n))。需要根据具体应用场景选择合适的存储结构。

	\textbf{【核心考点二：链表结点操作】}

	链表操作的核心是指针的修改。插入结点时需要正确处理指针的顺序，删除结点时需要记得释放内存。常见的操作包括：在指定结点后插入、删除指定结点、在表头插入删除、在表尾插入删除等。循环链表还需要注意表尾指针的指向。

	\textbf{【核心考点三：时间复杂度分析】}

	顺序表：按位置访问O(1)、按位置插入删除O(n)。链表：按位置访问O(n)、按位置插入删除O(1)。这些是选择存储结构的重要依据。需要特别注意“按位置访问”和“按位置插入删除”的区别。

	\textbf{【核心考点四：头结点的作用】}

	头结点是链表中的特殊结点，不存储实际数据。头结点的作用是简化链表操作，使得在表头的插入删除操作与在表中其他位置的插入删除操作统一处理。判断带头结点的单链表为空的条件是head->next == NULL。

	\textbf{【核心考点五：循环链表】}

	循环链表形成闭环，尾结点的指针指向头结点。循环单链表判断为空的条件是head == NULL（不带头结点）或head->next == head（带头结点）。循环双链表判断为空的条件是L->next == L。循环链表可以从任意结点出发遍历整个链表。
\end{bluebox}

\section{典型例题精讲}

\subsection{选择题精析}

\begin{bluebox}{答案与深度解析}
	\textbf{答案：A}

	\textbf{【核心认知框架】}

	本题考查链表的基本特点。链表是一种链式存储结构的线性表，其核心特点是通过指针连接各个结点。链表不支持随机访问，这是链表与顺序表最重要的区别之一。理解链表的访问特性是掌握链表操作的基础。

	\textbf{【选项原理深度拆解】}

	\item \textbf{A. 可随机访问任一结点 — 错误}
	\begin{itemize}
		\item 链表不支持随机访问，只能顺序访问
		\item 访问第i个结点需要从头开始遍历，时间复杂度为O(n)
		\item 这是链表与顺序表最核心的区别
	\end{itemize}

	\item \textbf{B. 插入删除不需要移动元素 — 正确}
	\begin{itemize}
		\item 链表插入删除只需要修改指针，不需要移动元素
		\item 时间复杂度为O(1)（前提是已经定位到位置）
		\item 这是链表相比顺序表最大的优势
	\end{itemize}

	\item \textbf{C. 不必事先估计存储空间 — 正确}
	\begin{itemize}
		\item 链表使用动态内存分配，可以随时创建新结点
		\item 不需要预先知道数据规模
		\item 内存按需分配，更加灵活
	\end{itemize}

	\item \textbf{D. 所需空间与其长度成正比 — 正确}
	\begin{itemize}
		\item 链表空间与元素个数成正比
		\item n个元素需要n个结点，每个结点包含数据和指针
		\item 空间复杂度为O(n)
	\end{itemize}

	\textbf{【应背诵知识点】}

	链表的核心特点：插入删除O(1)、访问O(n)、空间O(n)、存储密度<1。
\end{bluebox}

\begin{bluebox}{答案与深度解析}
	\textbf{答案：B}

	\textbf{【核心认知框架】}

	本题考查带头结点的单链表的空表判断条件。头结点是链表中一个特殊的结点，它不存储实际数据，仅作为链表的入口点。理解头结点的作用和空表的判定条件是链表操作的基础。

	\textbf{【选项原理深度拆解】}

	\item \textbf{A. head == NULL — 错误}
	\begin{itemize}
		\item 带头结点的单链表，头结点始终存在，head指针不为空
		\item head指向头结点，头结点是一个有效的结点
	\end{itemize}

	\item \textbf{B. head->next == NULL — 正确}
	\begin{itemize}
		\item 带头结点的单链表，头结点存在但后续没有其他结点
		\item 此时链表为空，只有头结点
	\end{itemize}

	\item \textbf{C. head->next == head — 错误}
	\begin{itemize}
		\item 这是循环链表的空表判定条件
		\item 循环链表中，头结点的next指向自己时表示为空
	\end{itemize}

	\item \textbf{D. head != NULL — 错误}
	\begin{itemize}
		\item 任何情况下，head指针都不为空
		\item 头结点始终存在，head指向头结点
	\end{itemize}

	\textbf{【应背诵知识点】}

	不带头结点的单链表为空的条件：head == NULL。带头结点的单链表为空的条件：head->next == NULL。循环单链表为空的条件：head == NULL（不带头结点）或head->next == head（带头结点）。
\end{bluebox}

\begin{bluebox}{答案与深度解析}
	\textbf{答案：D}

	\textbf{【核心认知框架】}

	本题考查线性表基本概念。线性表是具有n个数据元素的有限序列。n是线性表的长度，n≥0。当n=0时，线性表为空表。线性表可以是空表，这是一个重要的概念。

	\textbf{【选项原理深度拆解】}

	\item \textbf{A. 一个有限序列，可以为空 — 正确}
	\begin{itemize}
		\item 线性表是有限序列
		\item 长度为0时，线性表为空
		\item 符合线性表的定义
	\end{itemize}

	\item \textbf{B. 一个有限序列，不能为空 — 错误}
	\begin{itemize}
		\item 线性表可以为空表
		\item 这是错误选项
	\end{itemize}

	\item \textbf{C. 一个无限序列，可以为空 — 错误}
	\begin{itemize}
		\item 线性表是有限序列，不是无限序列
		\item 无限序列不是线性表
	\end{itemize}

	\item \textbf{D. 一个无限序列，不能为空 — 错误}
	\begin{itemize}
		\item 线性表是有限序列
		\item 无限序列不是线性表
	\end{itemize}

	\textbf{【应背诵知识点】}

	线性表是n（n≥0）个数据元素的有限序列。n=0时为空表。线性表不能是无限序列。
\end{bluebox}

\begin{bluebox}{答案与深度解析}
	\textbf{答案：C}

	\textbf{【核心认知框架】}

	本题考查顺序表与链表在不同操作上的效率对比。顺序表访问元素的时间复杂度为O(1)，链表访问元素的时间复杂度为O(n)。理解两种存储结构在不同操作上的效率差异是选择存储结构的依据。

	\textbf{【选项原理深度拆解】}

	\item \textbf{A. 输出第i个元素值 — 顺序表更高效}
	\begin{itemize}
		\item 顺序表可以直接通过下标访问，时间复杂度O(1)
		\item 链表需要遍历，时间复杂度O(n)
	\end{itemize}

	\item \textbf{B. 交换第0个元素与第1个元素的值 — 顺序表更高效}
	\begin{itemize}
		\item 两者时间复杂度相同，都是O(1)
		\item 顺序表直接交换数组元素
		\item 链表需要修改两个结点的数据域
	\end{itemize}

	\item \textbf{C. 顺序输出这n个元素的值 — 两者效率相当}
	\begin{itemize}
		\item 顺序表遍历O(n)
		\item 链表遍历O(n)
		\item 时间复杂度相同
	\end{itemize}

	\item \textbf{D. 输出与给定值x相等的元素序号 — 链表更高效}
	\begin{itemize}
		\item 两者平均时间复杂度都是O(n)
		\item 没有明显优势
	\end{itemize}

	\textbf{【答案解析】}

	本题要求找出在顺序表上实现比链表更高效的操作。选项A“输出第i个元素值”是典型的随机访问操作，顺序表可以直接通过下标计算地址访问，时间复杂度为O(1)；链表需要从头遍历，时间复杂度为O(n)。因此，顺序表在随机访问方面具有明显优势。

	\textbf{【应背诵知识点】}

	顺序表访问操作O(1)是核心优势。链表的优势在于插入删除O(1)。
\end{bluebox}

\begin{bluebox}{答案与深度解析}
	\textbf{答案：D}

	\textbf{【核心认知框架】}

	本题考查链表操作的时间复杂度分析。在已经定位到操作位置的情况下，链表的插入和删除操作只需要修改指针，时间复杂度为O(1)。这是链表相比顺序表最大的优势。

	\textbf{【选项原理深度拆解】}

	\item \textbf{A. 删除单链表中的第一个元素 — O(1)}
	\begin{itemize}
		\item 只需要修改头指针指向第二个结点
		\item 时间复杂度O(1)
	\end{itemize}

	\item \textbf{B. 删除单链表中的最后一个元素 — O(n)}
	\begin{itemize}
		\item 需要找到倒数第二个结点
		\item 需要遍历整个链表
		\item 时间复杂度O(n)
	\end{itemize}

	\item \textbf{C. 在单链表第一个元素前插入一个新元素 — O(1)}
	\begin{itemize}
		\item 在表头插入只需要修改头指针
		\item 时间复杂度O(1)
	\end{itemize}

	\item \textbf{D. 在单链表最后一个元素后插入一个新元素 — O(n)}
	\begin{itemize}
		\item 需要遍历到表尾
		\item 时间复杂度O(n)
	\end{itemize}

	\textbf{【答案解析】}

	在只有头指针的单链表中，删除最后一个元素需要找到倒数第二个结点，这需要遍历整个链表，时间复杂度为O(n)。而删除第一个元素、在第一个元素前插入、在最后一个元素后插入这些操作都可以在O(1)时间内完成。

	\textbf{【应背诵知识点】}

	单链表头指针操作：删除第一个元素O(1)、在第一个元素前插入O(1)。需要遍历的操作：删除最后一个元素O(n)、在最后一个元素后插入O(n)。
\end{bluebox}

\subsection{填空题精析}

\begin{enumerate}
	\item \textbf{在单链表中，要删除某一指定的结点，必须找到该结点的（前驱）结点。}

	      解析：单链表的结点只包含指向后继的指针，删除某个结点时需要修改其前驱结点的指针指向其后继结点。因此，必须先找到待删除结点的前驱结点。

	\item \textbf{在双链表中，每个结点有两个指针域，一个指向（前驱结点），另一个指向（后继结点）。}

	      解析：双链表的最大特点是每个结点有两个指针域，分别指向该结点的前驱和后继。这使得双链表可以双向遍历。

	\item \textbf{线性表常用的存储结构有（顺序存储）和（链式存储）。}

	      解析：线性表的两种基本存储方式是顺序存储（如顺序表）和链式存储（如链表）。这是线性表最基本的分类方式。

	\item \textbf{根据线性表的链式存储结构中每一个结点包含的指针个数，将线性链表分成（单链表）和（双链表）。}

	      解析：单链表每个结点包含一个指针（指向后继），双链表每个结点包含两个指针（分别指向前驱和后继）。

	\item \textbf{顺序存储结构是通过（下标）表示元素之间的关系的；链式存储结构是通过（指针）表示元素之间的关系的。}

	      解析：顺序存储利用数组的下标表示相邻关系，链式存储利用指针显式表示逻辑关系。这是两种存储方式的本质区别。

	\item \textbf{带头结点的循环链表L中只有一个元素结点的条件是（L->next->next == L）。}

	      解析：循环链表中，尾结点的next指向头结点。有一个元素结点时，头结点的next指向该元素，该元素的next指向头结点，即L->next->next == L。

	\item \textbf{在长度为n的顺序表中，删除第i个元素时，需向前移动（n-i）个元素。}

	      解析：删除第i个元素后，第i+1到第n个元素都需要向前移动一位，共移动n-i个元素。
\end{enumerate}

\subsection{判断题精析}

\begin{enumerate}
	\item \textbf{顺序存储的线性表可以随机存取。（正确）}

	      解析：顺序表使用数组实现，数组支持随机访问，可以通过下标直接访问任意位置的元素，时间复杂度为O(1)。

	\item \textbf{线性表中的元素可以是各种类型的，但同一线性表中的数据元素具有相同的性质，因此属于同一数据对象。（正确）}

	      解析：线性表中的元素必须属于同一数据对象，即具有相同的数据类型。这是线性表定义的基本要求。

	\item \textbf{在线性表的顺序存储结构中，逻辑上相邻的两个元素在物理位置上也是相邻的。（正确）}

	      解析：顺序存储的特点就是逻辑上相邻的元素在物理位置上也相邻。这是顺序存储的定义。

	\item \textbf{线性表的链式存储结构优于顺序存储结构。（错误）}

	      解析：链式存储和顺序存储各有优缺点，适用于不同场景，不存在绝对的优劣之分。选择哪种存储方式需要根据具体问题决定。

	\item \textbf{在单链表中，给定任一结点的地址p，则可用下述语句将新结点s插入结点p的后面：p->next = s; s->next = p->next; （错误）}

	      解析：这两条语句的顺序是错误的。正确的顺序应该是先设置s->next = p->next，再设置p->next = s。如果顺序颠倒，p->next的值已经被修改为s，无法找到原来的后继结点。

	\item \textbf{链表的每个结点中都恰好包含一个指针。（错误）}

	      解析：单链表的每个结点包含一个指针，双链表的每个结点包含两个指针。并非所有链表都只有一个指针。

	\item \textbf{双向链表可随机访问任一结点。（错误）}

	      解析：双向链表只是增加了前驱指针，便于找到前驱结点，但仍然不能随机访问。访问任意结点仍然需要从头遍历，时间复杂度为O(n)。
\end{enumerate}

\section{本章知识体系图}

\begin{center}
	\begin{tikzpicture}[scale=0.85, every node/.style={scale=0.85}]
		\tikzstyle{bx} = [rectangle, rounded corners, minimum width=2.5cm, minimum height=0.8cm, text centered, draw=black, fill=blue!20]
		\tikzstyle{ar} = [->, >=stealth, thick]

		\node[bx] (linear) at (0,0) {线性表};
		\node[bx] (seq) at (-3,-2) {顺序存储（顺序表）};
		\node[bx] (link) at (3,-2) {链式存储（链表）};
		\node[bx] (seqchar) at (-5,-4) {存储密度高};
		\node[bx] (seqacc) at (-3,-4) {访问O(1)};
		\node[bx] (seqins) at (-1,-4) {插入删除O(n)};
		\node[bx] (linkchar) at (1,-4) {存储密度低};
		\node[bx] (linkacc) at (3,-4) {访问O(n)};
		\node[bx] (linkins) at (5,-4) {插入删除O(1)};

		\node[bx] (single) at (1,-6) {单链表};
		\node[bx] (double) at (3.5,-6) {双链表};
		\node[bx] (circular) at (6,-6) {循环链表};

		\draw[ar] (linear) -- (seq);
		\draw[ar] (linear) -- (link);
		\draw[ar] (seq) -- (seqchar);
		\draw[ar] (seq) -- (seqacc);
		\draw[ar] (seq) -- (seqins);
		\draw[ar] (link) -- (linkchar);
		\draw[ar] (link) -- (linkacc);
		\draw[ar] (link) -- (linkins);
		\draw[ar] (link) -- (single);
		\draw[ar] (link) -- (double);
		\draw[ar] (link) -- (circular);
	\end{tikzpicture}
\end{center}

\section{本章学习建议}

本章是数据结构课程的重点章节之一，线性表是后续各章的基础。建议考生从以下方面深入学习：

第一，理解顺序存储和链式存储的原理。两种存储方式的核心区别在于元素之间的关系如何表示：顺序存储通过物理位置相邻表示，链式存储通过指针显式表示。

第二，熟练掌握各种链表操作的算法实现。重点是指针的修改顺序和边界条件的处理。建议考生自己动手编写代码，通过实际练习加深理解。

第三，重点掌握时间空间复杂度分析。顺序表和链表在不同操作上的时间复杂度是考试重点，需要熟练掌握并能够灵活应用。

第四，理解存储结构选择的原则。根据具体应用场景的特点（访问频繁还是插入删除频繁、数据规模是否固定等）选择合适的存储结构。

第五，多做练习，熟能生巧。链表操作是数据结构的基础技能，只有通过大量的练习才能真正掌握。

\end{document}